% Generated draft: requires mathematical and source review.
\documentclass[11pt]{article}
\usepackage[margin=25mm]{geometry}
\usepackage{fontspec}
\setmainfont{Latin Modern Roman}
\usepackage{amsmath,amssymb,mathtools,unicode-math}
\setmathfont{Latin Modern Math}
\usepackage{enumitem,needspace,xurl,hyperref}
\setlength{\parindent}{0pt}
\setlength{\parskip}{5pt}
\setlength{\emergencystretch}{3em}
\newcommand{\N}{\mathbb{N}}
\newcommand{\Z}{\mathbb{Z}}
\newcommand{\Q}{\mathbb{Q}}
\newcommand{\R}{\mathbb{R}}
\newcommand{\C}{\mathbb{C}}
\newcommand{\F}{\mathbb{F}}
\newcommand{\A}{\mathbb{A}}
\newcommand{\PP}{\mathbb{P}}
\newcommand{\E}{\mathbb{E}}
\newcommand{\eps}{\varepsilon}
\newcommand{\dd}{\,\mathrm{d}}
\begin{document}
\section*{UnsolvedMath ID: 30006349}
\textbf{Source problem code:} catalogue entry
\subsection{Definitions and mathematical statement}
Let \(p\) be an odd prime, and let \(\omega\) be a generator of the multiplicative group \(\F_p^\times\). Use the commutator convention \([x,y]=x^{-1}y^{-1}xy\). The presentations below are understood relative to groups of exponent dividing \(p\) and nilpotency class at most two: in addition to the displayed relations, impose \(x^p=1\) and \([[x,y],z]=1\) for all group elements \(x,y,z\). Thus they specify quotients of the free group on the five displayed generators by these ambient identities and the displayed relators. The supplied source describes the resulting groups as having exponent \(p\) and nilpotency class two.

Define
\[
\begin{aligned}
G_1=\langle a,b,c,d,e\mid{}&[c,b]=[d,b]=[e,c]=[e,d]=1,\\
 &[d,c]=[b,a],\quad [e,b]=[c,a],\\
 &[e,a]=[d,a]^\omega\rangle,
\end{aligned}
\]
and
\[
\begin{aligned}
G_2=\langle a,b,c,d,e\mid{}&[d,a]=[e,a]=[c,b]=[d,b]=[e,c]=1,\\
 &[d,c]=[b,a],\quad [e,b]=[c,a]\rangle.
\end{aligned}
\]
Here \([d,a]^\omega\) means \([d,a]^m\), where \(m\in\Z\) represents \(\omega\in\F_p\); the exponent identity makes this independent of the representative.

A group invariant is an assignment \(I\) to groups, with values in a specified set, such that \(H\cong K\) implies \(I(H)=I(K)\). For arbitrary odd \(p\) and a chosen generator \(\omega\), find such an invariant that is simpler than the automorphism-group order and establish
\[
I(G_1)\ne I(G_2),
\]
thereby proving \(G_1\not\cong G_2\). The source does not specify a mathematical criterion for ``simpler,'' so this comparison remains a qualitative requirement requiring explanation in any proposed answer.

The supplied dated curation states that automorphism-group order already distinguishes these groups and that no simpler invariant answering the extracted question had been established. These are supplied source assessments, not independently verified conclusions of this draft.
\subsection{Short English statement}
For each odd prime \(p\) and chosen generator \(\omega\in\F_p^\times\), seek an isomorphism invariant that distinguishes the two presented groups of exponent \(p\) and nilpotency class two and is simpler than computing their automorphism-group orders. The meaning of ``simpler'' is unspecified. The supplied open status and literature assessment require independent review.
\subsection{Sources}
\begin{itemize}
\item Local registry: \texttt{data/unsolvedmath/problems.json}, supplied record with permanent UnsolvedMath ID \(30006349\), title ``Simple Invariant Distinguishing Two Class-Two Exponent-\(p\) Groups,'' source problem code \texttt{catalogue entry}. The record belongs to the Oberwolfach Reports Open Problems collection and gives proposed year 2025; no individual proposer is supplied.
\item Supplied primary citation: \emph{Computational Group Theory} (2025), Oberwolfach Reports, DOI \texttt{10.4171/owr/2025/27}, pp.\ 1389--1450. \url{https://doi.org/10.4171/owr/2025/27}. The supplied background describes Oberwolfach Reports as documenting workshops and their open-problem discussions.
\item Supplied literature source: ``Oberwolfach Report 27/2025.'' \url{https://ems.press/content/serial-article-files/51860}.
\item Supplied literature source: Vaughan-Lee, \emph{Automorphisms of class-two groups of prime exponent}. \url{https://arxiv.org/abs/1501.00678}.
\item Literature assessment dated exactly \texttt{2026-08-22}: the supplied curation labels the problem open, states that automorphism-group size separates the groups, and reports that no simpler invariant answering the extracted question has been established. This is supplied dated curation derived from the supplied OWR material, not an independently verified literature assessment; its status and cited claims require independent verification.
\item The registry labels the statement \texttt{corrected\_verified} and reports that the official Oberwolfach report was checked, the question isolated from surrounding exposition, and notation repaired for self-contained publication. These are supplied provenance claims; no source visit or independent verification is asserted here. The retained original question asks whether there is a ``simpler invariant'' for non-isomorphism testing of these groups.
\end{itemize}
\end{document}
